\documentclass[12pt]{article}
\usepackage[utf8]{inputenc}
\usepackage[T1]{fontenc}
\usepackage{hyperref}
\title{Fyer Language Resources for K-12 Curriculum Development}
\author{Compiled by Claude Sonnet 5.5 (claude-sonnet-5-5) for Andr\'as Kornai}
\date{}
\begin{document}
\maketitle

\section{Introduction}
Fyer (Fier; ISO 639-3: fie) is a West Chadic (Afro-Asiatic) language of Plateau State, Nigeria. Wikipedia gives 26,000 speakers (2000); Joshua Project gives a Fyer population of about 54,000 (Ethnologue 2016), so figures are approximate. Joshua Project places the Fyer in Mangu LGA, with a Fyer district in Pankshin LGA. Very little was found.

\section{Primary Educational Resources}
None found.

\section{Linguistic and Reference Materials}
\begin{itemize}
\item Wikipedia, Fyer language: \url{https://en.wikipedia.org/wiki/Fyer_language}
\end{itemize}

\section{Religious and Community Resources}
\begin{itemize}
\item Joshua Project, Fyer in Nigeria (describes declining knowledge of culture and language): \url{https://joshuaproject.net/people_groups/11785}
\end{itemize}

\section{Digital and Audio Resources}
None found.

\section{Curriculum Development Context}
No orthography, primer or school programme was found in the results. Development would start from documentation and community consultation in Plateau State.
\end{document}